\section{Model kart i sterowanie talią}
\label{sec:model-kart-talii}
%TODO: mozliwym drogom rowoju - dodac link do rozdalu
Reprezentacja kart i mechanizm sterowania talią stanowią podstawową warstwę
środowiska. Elementy te zaprojektowałem niezależnie od zasad
Ultimate Texas Hold’em, dzięki czemu mogą być wykorzystywane także przez
inne warianty pokera, o czym wspominam w rozdziale poświęconym możliwym drogom rozwoju.

\subsection{Reprezentacja karty}
\label{subsec:reprezentacja-karty}

Pojedynczą kartę zaimplementowałem jako niemodyfikowalny obiekt
\texttt{Card}, zawierający dwa pola: rangę oraz kolor. Rangi zakodowałem za
pomocą typu wyliczeniowego \texttt{Rank}, natomiast kolory za pomocą typu
\texttt{Suit}. Przez zastosowanie typów wyliczeniowych ograniczyłem możliwość
utworzenia karty zawierającej niepoprawną wartość, a przez podejście do karty jako niemutowalnego
obiektu zapewniłem bezpieczne umieszczanie jej w krotkach i zbiorach, a także możliwość wykorzystania
porównania liczebności zbioru do wykrywania zduplikowanych kart.

Rangi kart przyjmują wartości liczbowe od 2 do 14. Karty od dwójki do
dziesiątki zachowują swoje naturalne wartości, natomiast walet, dama, król
i as otrzymują odpowiednio wartości 11, 12, 13 i 14. Zdecydowałem się na taką konwencję 
ponieważ umożliwia mi to bezpośrednie porównywanie rang podczas oceny układów.

Warto dodać, że as jest kartą najwyższą, natomiast szczególny
przypadek strita od asa do piątki obsługiwany jest przez ewaluator układów. 
Zbiór kolorów obejmuje trefle, karo, kiery i piki.


\subsection{Standardowa talia}
\label{subsec:standardowa-talia}

Klasa \texttt{Deck} reprezentuje standardową talię. Podczas jej tworzenia
generowane są wszystkie kombinacje trzynastu rang i czterech kolorów.
Liczba kart w pełnej talii wynosi zatem

\begin{equation}
    |\mathcal{D}| = 13 \cdot 4 = 52.
\end{equation}

Ponieważ każda kombinacja rangi i koloru występuje dokładnie raz,
wygenerowana talia zawiera 52 unikalne karty.

Podstawową operacją jest metoda \texttt{draw()}, która pobiera jedną kartę
z wierzchu talii i następnie ją usuwa z kolekcji. Jest również
operacja służąca do pobrania wielu kart \texttt{draw\_many()}. Próba pobrania karty z pustej talii lub większej
liczby kart niż liczba dostępnych kart pozostałych w talii powoduje zgłoszenie wyjątku.
Chciałem, aby takie zachowanie umożliwiało szybkie wykrycie niepoprawnego przebiegu rozdania.

Domyślnie talia jest tasowana bezpośrednio po utworzeniu. Zadbałem o możliwość
utworzenia talii w konkretnej kolejności początkowej, co bywa użyteczne podczas testowania
samego modelu kart.

Aby zapewnić powtarzalność eksperymentów, w klasie \texttt{Deck} zastosowałem możliwość
przekazania ziarna generatora liczb pseudolosowych. Utworzenie dwóch talii z tym samym
ziarnem prowadzi do uzyskania tej samej kolejności kart. Na poziomie klasy \texttt{UTHGame}
zastosowałem dodatkowo generator nadrzędny, który przy każdym rozpoczęciu nowego rozdania
generuje osobne ziarno przekazywane do nowej talii. Dzięki temu kolejne rozdania w obrębie
jednego eksperymentu różnią się od siebie, ale cała ich sekwencja może zostać odtworzona przez
ponowne uruchomienie środowiska z tym samym ziarnem początkowym. Jeżeli dwa eksperymenty
zostaną uruchomione z tym samym ziarnem środowiska i zachowają tę samą kolejność wywołań,
otrzymają identyczne talie, co ogranicza wpływ losowego doboru kart na ocenę różnic
pomiędzy badanymi strategiami.

\subsection{Abstrakcja źródła kart i talia testowa}
\label{subsec:zrodlo-kart}

Silnik rozdania nie zależy bezpośrednio od klasy \texttt{Deck}. Zamiast tego
korzysta z abstrakcji \texttt{CardSource}, która wymaga jedynie udostępnienia
operacji \texttt{draw()}. Każdy obiekt spełniający ten kontrakt może zostać
użyty jako źródło kolejnych kart, dzięki czemu moduł odpowiedzialny za
rozdawanie nie musi rozpoznawać, czy korzysta z przetasowanej talii, czy
ze źródła deterministycznego.

Dzięki deterministycznemu źródłu byłem w stanie przygotować powtarzalne testy obejmujące
określone scenariusze rozgrywki, kwalifikacji krupiera i rozliczenia
zakładów. Do testowania konkretnych scenariuszy wprowadziłem klasę \texttt{FixedDeck},
która jest jawnie przekazaną sekwencją kart.
Przed zapisaniem sekwencji sprawdzane jest, czy wszystkie
jej elementy są poprawnymi obiektami \texttt{Card} oraz czy każda karta jest unikalna, 
dzięki czemu błąd w danych testowych nie może
doprowadzić do rozdania zawierającego dwie identyczne karty. 

\section{Ocena układów pokerowych}
\label{sec:ocena-ukladow}

Zadaniem modułu oceny układów pokerowych jest określenie wartości dowolnego
poprawnego układu pięciokartowego oraz wybranie najlepszej możliwej ręki
z pięciu, sześciu lub siedmiu dostępnych kart. W finalnej fazie podczas showdownu
ewaluator jest wywoływany dla gracza i krupiera zawsze na siedmiu kartach. 
Pozostałe dwa warianty wykorzystuje natomiast agent regułowy
opisany w rozdziale poświęconym agentom bazowym.

\subsection{Reprezentacja wartości ręki}
\label{subsec:wartosc-reki}

Do reprezentacji kategorii układów pokerowych podszedłem implementując typ wyliczeniowy \texttt{HandRank}.
Kolejnym kategoriom przypisuje rosnące wartości liczbowe od 0 do 8, począwszy od wysokiej karty,
a kończąc na pokerze. Dzięki temu porównanie kategorii odbywa się wprost 
przez porównanie odpowiadających im wartości całkowitych.

Sama kategoria nie wystarcza jednak do rozstrzygnięcia wszystkich przypadków (np. gracz i krupier mogą mieć układy tej samej kategorii, lecz różniące się rangą kart
tworzących dany układ albo kartami dodatkowymi). Z tego względu pełną wartość ręki
zaimplementowałem jako obiekt \texttt{HandValue}, zawierający kategorię
układu \texttt{rank} oraz uporządkowaną krotkę wartości rozstrzygających
remis \texttt{tiebreak}.

Klucz porównawczy ręki można zapisać jako

\begin{equation}
    K(h) =
    \left(
        r(h),
        t_1(h),
        t_2(h),
        \ldots,
        t_n(h)
    \right),
\end{equation}

gdzie \(r(h)\) oznacza wartość kategorii układu, natomiast
\(t_1(h),\ldots,t_n(h)\) są kolejnymi wartościami rozstrzygającymi remis.
Klucze porównywane są leksykograficznie. W pierwszej kolejności
porównywana jest kategoria układu, a dopiero w przypadku jej równości 
kolejne elementy krotki \texttt{tiebreak}.

Dla lepszego zobrazowania, wartość pary asów z królem, dziesiątką i dziewiątką jako kartami
dodatkowymi może zostać zapisana w postaci

\begin{equation}
    K(h) =
    \left(
        \text{ONE\_PAIR},
        14,
        13,
        10,
        9
    \right).
\end{equation}

Jeżeli druga ręka również zawiera parę asów, porównanie przechodzi kolejno
do króla, dziesiątki i dziewiątki. Kolory kart nie uczestniczą w rozstrzyganiu
remisów, zgodnie ze standardowymi zasadami pokera.

Zadbałem, aby obiekt \texttt{HandValue} sprawdzał zgodność długości krotki
\texttt{tiebreak} z kategorią układu, walidował również, czy wszystkie
jej elementy są liczbami całkowitymi odpowiadającymi poprawnym rangom kart.
W ten sposób ograniczyłem możliwość utworzenia wartości ręki, która nie mogłaby powstać
w prawidłowym rozdaniu.

\subsection{Ocena układu pięciokartowego}
\label{subsec:ocena-pieciu-kart}

Funkcja \texttt{evaluate\_five\_card\_hand()} przyjmuje pięć
unikalnych kart. W pierwszym kroku funkcja oblicza liczbę wystąpień każdej rangi. Na tej
podstawie możliwe jest wykrycie par, dwóch par, trójki, fulla oraz karety.
Niezależnie sprawdzane są dwa warunki: czy wszystkie karty mają ten sam
kolor oraz czy pięć różnych rang tworzy ciąg kolejnych wartości.

Kategorie sprawdzane są od najsilniejszej do najsłabszej. Jeżeli spełnione
są jednocześnie warunki strita i koloru, zwracany jest poker. W przeciwnym
razie funkcja sprawdza kolejno kareta, full, kolor, strit, trójka, dwie pary,
jedna para oraz wysoka karta.

Wartość \texttt{tiebreak} jest budowana bezpośrednio podczas rozpoznawania
kategorii. Wynikiem funkcji nie jest jedynie nazwa układu, lecz kompletna
wartość umożliwiająca porównanie z dowolną inną poprawną ręką. Szczególnym
przypadkiem jest strit \(A\)-\(2\)-\(3\)-\(4\)-\(5\), w którym as pełni rolę
najniższej karty, ewaluator zwraca wtedy strita o najwyższej karcie równej 5,
bez zmiany wartości asa w pozostałych kategoriach.

\subsection{Wybór najlepszej ręki}
\label{subsec:wybor-najlepszej-reki}

Podczas showdownu gracz i krupier dysponują łącznie siedmioma kartami. 
Rękę pokerową tworzy dokładnie pięć kart, dlatego funkcja \texttt{evaluate\_best\_hand()}
generuje wszystkie możliwe pięciokartowe podzbiory dostępnych kart, ocenia
każdy z nich, a następnie wybiera układ o największej wartości \texttt{HandValue}.
Generowanie każdego możliwego podzbioru nie jest optymalnym rozwiązaniem,
jednak liczba kombinacji jest niewielka i stała, co potraktowałem jako kompromis pomiędzy prostotą implementacji a wydajnością.

Liczba sprawdzanych kombinacji dla \(n\) dostępnych kart wynosi

\begin{equation}
    N(n) = \binom{n}{5}.
\end{equation}

Dla liczby kart obsługiwanych przez ewaluator jest tyle kombinacji:

\begin{equation}
    \binom{5}{5} = 1,
    \qquad
    \binom{6}{5} = 6,
    \qquad
    \binom{7}{5} = 21.
\end{equation}

W przypadku standardowego showdownu UTH konieczne jest zatem sprawdzenie
21 kombinacji dla gracza oraz 21 kombinacji dla krupiera.

Przed wygenerowaniem kombinacji sprawdzane jest, czy przekazano od pięciu
do siedmiu kart, czy wszystkie elementy są kartami oraz czy żadna karta
nie występuje więcej niż raz.

Wynikiem operacji jest obiekt \texttt{EvaluatedHand}, który przechowuje
zarówno wartość \texttt{HandValue}, jak i dokładnie pięć kart tworzących
wybrany układ. Uznałem to za wygodny sposób prezentacji wyniku showdownu oraz diagnostyki evaluatora.

Dla najwyższej kategorii, jaką jest poker królewski, nie tworzyłem osobnej kategorii \texttt{HandRank}.
Rozwinąłem go jako poker z asem jako najwyższą kartą.

